\subsection{Testing}
\label{sec:impl_testing}

This section presents the test execution results for the implemented modules of the first iteration. The backend test suite uses \path{pytest} with \texttt{asyncio\_mode = "auto"} for native async test support, running against a real PostgreSQL database in integration mode. Tests are organised into unit tests (pure function isolation) and integration tests (full request--response cycles).

Table~\ref{tab:test_result_summary} summarises the test execution results for the modules delivered in this iteration.

\begin{table}[htbp]
\centering
\caption{Test Execution Summary by Module (Iteration 1)}
\label{tab:test_result_summary}
\renewcommand{\arraystretch}{1.3}
\begin{tabular}{|l|c|c|c|}
\hline
\textbf{Module} & \textbf{Test Cases} & \textbf{Passed} & \textbf{Result} \\
\hline
Authentication \& Access Control & 8 & 8 & \textbf{Pass} \\
Course \& Material Management    & 5 & 5 & \textbf{Pass} \\
Quiz Generation \& Taking        & 4 & 4 & \textbf{Pass} \\
\hline
\textbf{Total}                   & \textbf{17} & \textbf{17} & \textbf{100\%} \\
\hline
\end{tabular}
\end{table}

% ------------------------------------------------------------------
\subsubsection{Authentication and Access Control}
\label{sec:impl_test_auth}

Table~\ref{tab:tc_detail_auth} details the test cases covering the authentication and access control module.

\renewcommand{\arraystretch}{1.25}
\small
\begin{longtable}{|>{\centering\arraybackslash}m{0.10\textwidth}|>{\raggedright\arraybackslash}m{0.30\textwidth}|>{\raggedright\arraybackslash}m{0.35\textwidth}|>{\centering\arraybackslash}m{0.08\textwidth}|}
\caption{Test Results: Authentication and Access Control}
\label{tab:tc_detail_auth} \\
\hline
\textbf{TC-ID} & \textbf{Description} & \textbf{Test File} & \textbf{Result} \\ \hline
\endfirsthead
\multicolumn{4}{c}{{\bfseries Table \thetable\ -- Continued}} \\ \hline
\textbf{TC-ID} & \textbf{Description} & \textbf{Test File} & \textbf{Result} \\ \hline
\endhead
\hline
\endfoot
\hline
\endlastfoot

TC-1.1.1 & Google SSO login callback & \path{test_oauth_flow.py} & Pass \\ \hline
TC-1.1.2 & Idempotent logout & \path{test_oauth_flow.py} & Pass \\ \hline
TC-1.2.1 & Session creation and refresh-token rotation & \path{test_identity_full.py} & Pass \\ \hline
TC-1.2.2 & Session revocation on account disable & \path{test_identity_queries.py} & Pass \\ \hline
TC-1.3.1 & TOTP verification and recovery-code consumption & \path{test_mfa_router.py} & Pass \\ \hline
TC-1.3.2 & Recovery-code reuse prevention & \path{test_mfa_gate.py} & Pass \\ \hline
TC-1.5.1 & Recursive permission resolution & \path{test_access_control_full.py} & Pass \\ \hline
TC-1.5.2 & Endpoint enforcement by role & \path{test_access_control_membership_resolution.py} & Pass \\ \hline

\end{longtable}
\normalsize

% ------------------------------------------------------------------
\subsubsection{Course and Material Management}
\label{sec:impl_test_course}

Table~\ref{tab:tc_detail_course} details the test cases for course and material management features implemented in this iteration.

\renewcommand{\arraystretch}{1.25}
\small
\begin{longtable}{|>{\centering\arraybackslash}m{0.10\textwidth}|>{\raggedright\arraybackslash}m{0.30\textwidth}|>{\raggedright\arraybackslash}m{0.35\textwidth}|>{\centering\arraybackslash}m{0.08\textwidth}|}
\caption{Test Results: Course and Material Management}
\label{tab:tc_detail_course} \\
\hline
\textbf{TC-ID} & \textbf{Description} & \textbf{Test File} & \textbf{Result} \\ \hline
\endfirsthead
\multicolumn{4}{c}{{\bfseries Table \thetable\ -- Continued}} \\ \hline
\textbf{TC-ID} & \textbf{Description} & \textbf{Test File} & \textbf{Result} \\ \hline
\endhead
\hline
\endfoot
\hline
\endlastfoot

TC-3.1.1 & Course lifecycle (draft $\to$ publish) & \path{test_courses_e2e.py} & Pass \\ \hline
TC-3.1.2 & Course archive and restore & \path{test_courses_e2e.py} & Pass \\ \hline
TC-3.3.1 & PDF material upload and ingestion trigger & \path{test_materials_e2e.py} & Pass \\ \hline
TC-3.3.2 & Ingestion pipeline processing & \path{test_materials_worker.py} & Pass \\ \hline
TC-3.7.1 & Student material access (presigned URL) & \path{test_materials_learner_router.py} & Pass \\ \hline

\end{longtable}
\normalsize

% ------------------------------------------------------------------
\subsubsection{Quiz Generation and Taking}
\label{sec:impl_test_quiz}

Table~\ref{tab:tc_detail_quiz} details the test cases for the AI quiz generation and student quiz-taking features.

\renewcommand{\arraystretch}{1.25}
\small
\begin{longtable}{|>{\centering\arraybackslash}m{0.10\textwidth}|>{\raggedright\arraybackslash}m{0.30\textwidth}|>{\raggedright\arraybackslash}m{0.35\textwidth}|>{\centering\arraybackslash}m{0.08\textwidth}|}
\caption{Test Results: Quiz Generation and Taking}
\label{tab:tc_detail_quiz} \\
\hline
\textbf{TC-ID} & \textbf{Description} & \textbf{Test File} & \textbf{Result} \\ \hline
\endfirsthead
\multicolumn{4}{c}{{\bfseries Table \thetable\ -- Continued}} \\ \hline
\textbf{TC-ID} & \textbf{Description} & \textbf{Test File} & \textbf{Result} \\ \hline
\endhead
\hline
\endfoot
\hline
\endlastfoot

TC-4.1.1 & AI quiz generation from PDF content & \path{test_quiz_full_pipeline.py} & Pass \\ \hline
TC-4.1.2 & Teacher review (approve/reject/edit) & \path{test_quiz_authoring_service.py} & Pass \\ \hline
TC-4.3.1 & Student quiz session submission & \path{test_quiz_taking_service.py} & Pass \\ \hline
TC-4.3.2 & Per-question response and feedback & \path{test_quizzes_e2e.py} & Pass \\ \hline

\end{longtable}
\normalsize

% ------------------------------------------------------------------
\subsubsection{CI Pipeline}
\label{sec:impl_ci}

The GitHub Actions workflow (\texttt{.github/workflows/backend-new-ci.yml}) runs the following jobs on every pull request:

\begin{enumerate}
    \item \textbf{Lint} — \texttt{ruff check}, \texttt{ruff format --check}, \texttt{mypy abridgeai}, \texttt{bandit -r abridgeai/}.
    \item \textbf{Secret scan} — \texttt{gitleaks detect} with the project's custom ruleset covering OpenAI, Anthropic, and Garage S3 key patterns.
    \item \textbf{Unit tests} — \path{pytest tests/unit/} with coverage reporting.
    \item \textbf{Integration tests} — \path{pytest tests/integration/} against a PostgreSQL service container spun up by the workflow.
\end{enumerate}

All pull requests must pass all four stages before merge is permitted. Live infrastructure tests (S3, Neo4j, audio, video) are gated behind pytest markers and excluded from the default CI run to keep pipeline duration manageable.
